% =============================================================================
% Chapitre 10 : Limites et perspectives
% =============================================================================
\chapter{Limites et perspectives}
\label{ch:limitations}

Ce chapitre évalue honnêtement les limites actuelles du système et identifie des directions d'amélioration concrètes.

\section{Limites actuelles}

\subsection{Limites liées aux données}

\textbf{Jeu spam réduit et hors domaine} : la SMS Spam Collection ($\sim$5\,500 lignes) est petite face au jeu de toxicité ($\sim$160K), et le vocabulaire propre aux SMS (« free », « winner », « claim ») ne transfère probablement ni au spam web, ni au spam e-mail ou réseaux sociaux. Le modèle n'a pas appris de pattern de détection de spam significatif dans les expériences sur sous-ensemble.

\textbf{Couverture du discours haineux} : les étiquettes haine proviennent uniquement des colonnes \texttt{identity\_hate} et \texttt{threat} du jeu Jigsaw, captant une tranche étroite du phénomène---sans haine implicite, langage codé ni discrimination contextuelle.

\textbf{Problème des étiquettes partielles} : le masquage ($y_k = -1$) est une solution pratique aux données hétérogènes, mais le modèle ne reçoit aucun signal de supervision sur les relations inter-catégories (ex. il ne peut pas apprendre qu'un message spam est rarement du discours haineux, ces étiquettes étant toujours masquées pour le jeu spam).

\textbf{Pas de données image étiquetées} : le modèle image n'a jamais été évalué contre un ensemble de test étiqueté ; la précision zéro-shot de CLIP reste inconnue pour cette configuration précise de prompts et de seuils.

\subsection{Limites du modèle}

\textbf{Contrainte CPU} : l'entraînement de BERT sur le jeu complet sur CPU est prohibitif ($\sim$22 s par lot), forçant l'usage de petits sous-ensembles et de peu d'époques, d'où les résultats sous-optimaux du chapitre~\ref{ch:results}.

\textbf{Capacité fixe} : BERT-base (110M paramètres) est limité face aux modèles plus grands (BERT-large, RoBERTa-large), et aucun pré-entraînement spécifique au domaine de la modération n'est exploité.

\textbf{Scores image non calibrés} : la pseudo-probabilité NSFW de CLIP n'est pas calibrée ; les seuils (0,30, 0,70) sont heuristiques et exigent une calibration sur données de validation.

\subsection{Limites du système}

\textbf{Pas d'apprentissage temps réel} : le réentraînement est hebdomadaire ; les erreurs découvertes entre deux cycles ne sont pas corrigées immédiatement. \textbf{Langue unique} : \texttt{bert-base-uncased} ne supporte que l'anglais. \textbf{Pas de vidéo ni audio} : le système ne traite que texte et images. \textbf{Authentification limitée} : l'infrastructure JWT (\texttt{python-jose}, \texttt{passlib}) existe mais le middleware d'authentification n'est pas implémenté.

\section{Perspectives}

\textbf{Priorité haute} :
\begin{enumerate}[nosep]
  \item Entraînement complet sur les 164K lignes avec 4+ époques dans un environnement GPU---le chemin le plus direct vers une modération texte de qualité production.
  \item Jeu d'évaluation image étiqueté pour les cinq catégories, mesure des performances zéro-shot et calibrage des seuils de routage.
  \item Apprentissage actif sélectionnant les images les plus incertaines pour la revue humaine, accélérant l'accumulation de données pour la sonde linéaire.
  \item Remplacement ou complément du jeu SMS par des jeux spam plus grands et mieux adaptés au domaine (Enron Spam, SpamAssassin).
  \item Entraînement accéléré GPU (CUDA/MPS) pour itérer vite et exploiter le jeu complet.
\end{enumerate}

\textbf{Priorité moyenne} :
\begin{enumerate}[nosep]
  \item Modération multilingue via BERT multilingue ou XLM-RoBERTa.
  \item Calibration des probabilités (temperature scaling, Platt scaling) pour un choix de seuils fondé sur les taux de faux positifs/négatifs souhaités.
  \item Jeux dédiés au discours haineux (HateXplain, Implicit Hate Speech Corpus) pour détecter la haine implicite et codée.
  \item Routage sensible aux coûts, pondérant la décision par le coût relatif des faux positifs et faux négatifs par catégorie.
  \item Finalisation du middleware d'authentification JWT.
\end{enumerate}

\textbf{Priorité basse} :
\begin{enumerate}[nosep]
  \item Modération vidéo par échantillonnage d'images clés appliquant le pipeline image.
  \item Robustesse adversaire du modèle texte (substitutions de caractères, astuces Unicode, paraphrases).
  \item Distillation du modèle BERT vers un modèle plus petit et rapide pour l'embarqué ou le haut débit.
  \item Infrastructure de tests A/B de versions de modèles et configurations de routage.
  \item Confidentialité différentielle pour protéger les exemples d'entraînement individuels.
\end{enumerate}

\section{Considérations éthiques}

\begin{itemize}[nosep]
  \item \textbf{Biais des données} : le jeu Jigsaw présente des biais connus, dont la sur-signalisation de l'anglais afro-américain ; le système actuel hérite de ces biais.
  \item \textbf{Transparence} : le système fournit des motifs de routage lisibles, mais les décisions du modèle sous-jacent ne sont pas interprétables par nature.
  \item \textbf{Bien-être des modérateurs} : les relecteurs humains sont exposés à des contenus potentiellement nuisibles ; avertissements de contenu et politiques de rotation sont nécessaires.
  \item \textbf{Risque de sur-censure} : la faible précision toxique suggère un risque de censure excessive de contenus légitimes, particulièrement de communautés marginalisées.
\end{itemize}
